\documentclass{article}
\usepackage{graphicx}
\usepackage{listings}
\usepackage{tikz-cd}
\usepackage{tikz}
\usepackage{amsmath}
\usepackage{amssymb}
\usepackage{amsfonts}
\usepackage{wasysym}
\usepackage[brazilian]{babel}
\usepackage{hyperref}
\usepackage{eso-pic,graphicx}
\usepackage{float}
\usepackage{pgfplots}
\usepackage{caption}
\captionsetup{font=footnotesize}
\usepackage{multirow}
\usepackage{xcolor}
\usepackage{bold-extra}
\usepackage[margin=2cm]{geometry}
\usepackage{titlesec}
\usepackage{adjustbox} % Para redimensionar a tabela

% Define a formatação das seções
\titleformat{\section}
  {\normalfont\large\bfseries} % Formato do título
  {\thesection} % Número da seção
  {1em} % Espaçamento entre número e título
  {} % Texto antes do título

\titleformat{\subsection}
  {\normalfont\small\bfseries} % Formato do título
  {\thesubsection} % Número da subseção
  {1em} % Espaçamento entre número e título
  {} % Texto antes do título

\titleformat{\subsubsection}
  {\normalfont\small\bfseries} % Formato do título
  {\thesubsubsection} % Número da subsubseção
  {1em} % Espaçamento entre número e título
  {} % Texto antes do título

\definecolor{minhacor}{RGB}{0, 128, 255} % Azul

%========================================================================================================================================

\begin{document}

	\begin{center}
		{\Large {Trabalho Prático I - Redes de Computadores}}
		
		\vspace{.2cm}
		
		{Gabriel de Araújo Costa Gomes}

        {2022102872}
		
		\vspace{.2cm}
		
		Departamento de Ciência da Computação; Universidade Federal de Minas Gerais
		
		Belo Horizonte, Minas Gerais, Brasil; {\today}

        \texttt{gacg2004@ufmg.br}
		
	\end{center}

    \section{Introdução}

        \hspace{0.5cm}Esse relátório tem o objetivo de descrever os problemas, imprevistos, desafios e a lógica por trás das soluções do primeiro trabalho prático
        da matéria de Redes de Computadores no semestre 2026/1. O objetivo do trabalho foi desenvolver um sistema de quebra senhas
        baseado na mecânica de tentativa e erro com \textit{feedback} posicional, seguindo um modelo cliente-servidor.

    \section{Desafios, Imprevistos, Dificuldades e Soluções}
        \subsection{Desafios}
        A partir da leitura do enunciado do trabalho já foram detectados os seguintes desafios:

        \begin{enumerate}
            \item Estabelecimento correto das conexões;
            \item Flexibilidade no uso dos domínios IPv4 ou IPv6;
            \item Tratamento e verificação de todas as entradas feitas ao sistema para ter o funcionamento desejado sempre;
            \item Tratamento das mensagens trocadas entre o cliente e o servidor; e por fim
            \item Tratamento das mensagens de feedback do server.
        \end{enumerate}

        \subsection{Fluxo de desenvolvimento}
        
        \hspace{0.5cm}Primeiro foi estudado o funcionamento das conexões de socket e como elas trocam mensagens
        entre si, em seguida foi feito o desenvolvimento das partes funcionais: a inicialização do server e do client. Depois de garantir que estava tudo funcional
        na conexão, foi detectado os lugares onde poderiam ser usadas funções para substituir, como por exemplo a função  \textit{ip\_validation} e \textit{server\_especifications}
        (esta que existe tanto no client quanto no server). Essa etapa foi puramente gosto pessoal, com o intuito de deixar mais legível o código e tornar mais simples de debbugar.
        
        
        Para a parte da comunicação, tudo iniciou com pequenos testes com o envio de mensagens simples, seguido do desenvolvimento das funções que fazem o comportamento de comunicação solicitado 
        pelo enunciado, sendo estas as funções \textit{guesses} no client e \textit{game} no server. 

        \subsection{Imprevistos, Dificuldades e soluções}

        \hspace{0.5cm}Durante o desenvolvimento que foi descrito, houveram alguns problemas na parte da comunicação. Um deles foi a validação de entrada de senha feita pelo cliente na função \textit{guesses},
        que necessitou de conversão de \textit{char} para \textit{int}. Como a função \textit{atoi} não iria me ajudar muito bem nesse caso, pois tiraria a separação dos números, foi utilizado a operação x - '0',
        onde x é um \textit{char} de um número qualquer. O restante da função foi bem tranquila.

        Já função de comunicação \textit{game} no server foi mais trabalhosa. A dificuldade foi achar uma forma de comparar dois vetores da forma necessária e formatar a saída de acordo. Para resolver isso, foi utilizado
        um vetor que mapeia a quantidade de dígitos de 0 a 9 existem na senha, o que facilitou bastante, mas pra que o feedback ficasse da forma correta foi muita tentativa e erro. O tratamento do feedback acabou tendo que 
        ser feito por etapas, pra evitar erros e melhorar legibilidade.

        No momento de utilizar essas funções aconteceram muitos problemas: mensagens que não chegavam, caracteres fantasmas, fechamento abrubto do server. Como não eram erros de conexão, pois isso já estava correto, tomei medidas
        para resolver esses caracteres estranhos com um \textit{memset} pra cada \textit{struct HackerMessage} em ambas as funções. Isso não resolveu o problema, então pensei que o problema pudesse ser a mensagem de feedback, então
        foi feito \textit{$feedback\_msg[5] = '\backslash 0'$}, também sem sucesso. Além disso, foi debbugado as funções de \textit{bind}, \textit{listen} e \textit{accept}, mas não houve erro algum nelas. No final de tudo, descobriu-se que o problema era a forma com que foi inicializada a função criada: com o socket do servidor. Um erro inocente mas que acredito que 
        foi muito importante para o aprendizado: os sockets do server não transferem dados, eles apenas criam a conexão e mantem a comunicação com o \textit{listen}. Depois de corrigir o socket, tudo funcionou normalmente, haviam uns erros
        no feedback ainda, mas eles foram corrigidos em seguida.

    \section{Conclusão}

    \hspace{0.5cm}O trabalho se mostrou bem sucedido em ensinar o funcionamento de uma conexão TCP/IP simples, reforçando os conceitos básicos tanto do lado do servidor quanto do cliente. Além disso, a atividade proporcionou uma compreensão mais 
    prática sobre a comunicação em rede, especialmente no envio e recebimento de dados, tratamento de erros e sincronização entre processos. 

    A proposta foi, ao mesmo tempo, desafiadora e divertida, e permitiu uma experiência profundada por colocar em prática os conceitos estudados. Durante o desenvolvimento, também foi possível perceber a importância de detalhes de 
    implementação, como o correto uso de estruturas, buffers e validação de dados, que impactam diretamente no funcionamento do sistema.

\end{document}